% Editorial bootstrap from the immutable source revision.
% This file is canonical and may be edited independently.

\part{Exact Inverse Moduli and Computability}


% Source fragment: Inverse computability; prefix co.
\section{Main result and proof architecture}

\subsection{The theorem in one statement}

Throughout the report, $F:\RealNumbers\to[0,1]$ denotes the bounded Fabius function:
it is $0$ on $(-\infty,0]$, $1$ on $[1,\infty)$, and on $[0,1]$ it is the
strictly increasing $C^\infty$ function satisfying
\begin{equation}\label{co:eq:F-basic}
 F(x)+F(1-x)=1,
 \qquad
 F'(x)=2F(2x)\quad(0\le x\le\tfrac12).
\end{equation}
Let $G:[0,1]\to[0,1]$ be its inverse, and define the totalized inverse
\begin{equation}\label{co:eq:totalized-inverse}
 \InvT(y):=G(\Clamp(y)),
 \qquad
 \Clamp(y):=\min(1,\max(0,y)).
\end{equation}
Thus $\InvT(y)=0$ for $y\le0$ and $\InvT(y)=1$ for $y\ge1$.

\begin{theorem}[Computability of the inverse Fabius function]\label{co:thm:main}
The inverse $G:[0,1]\to[0,1]$, in the subset-domain version of the
Grzegorczyk definition, and its totalization $\InvT:\RealNumbers\to\RealNumbers$ are computable
real functions.  More explicitly:
\begin{enumerate}
\item $G$ and $\InvT$ are sequentially computable.  From one algorithm that
      names a computable sequence $(y_i)$, one can uniformly compute dyadic
      approximations to $G(y_i)$ (or $\InvT(y_i)$) of any requested precision.
\item For $r\ge0$, put
      \begin{equation}\label{co:eq:Delta-intro}
       \Delta_r=\frac{1}{2^{r(r+1)/2}(r+1)^{r+1}}.
      \end{equation}
      Then, for all $u,v\in\RealNumbers$,
      \begin{equation}\label{co:eq:explicit-dyadic-modulus-intro}
       \AbsoluteValue{u-v}<\Delta_r
       \quad\Longrightarrow\quad
       \AbsoluteValue{\InvT(u)-\InvT(v)}<2^{-r}.
      \end{equation}
\item For $n\ge1$, let
      \begin{equation}\label{co:eq:r-of-n}
       r(n):=\min\SetOf{r\ge1:2^r>n},
      \end{equation}
      and define
      \begin{equation}\label{co:eq:d-of-n-intro}
       d(n):=2^{r(n)(r(n)+1)/2}\bigl(r(n)+1\bigr)^{r(n)+1},
       \qquad d(0):=1.
      \end{equation}
      The integer function $d$ is recursive (indeed primitive recursive after
      a standard bounded-search coding), and
      \begin{equation}\label{co:eq:wikipedia-modulus-intro}
       \AbsoluteValue{u-v}<\frac1{d(n)}
       \quad\Longrightarrow\quad
       \AbsoluteValue{\InvT(u)-\InvT(v)}<\frac1n
       \qquad(n\ge1).
      \end{equation}
\end{enumerate}
\end{theorem}

\begin{proof}
The claims are assembled below from independently certified ingredients; we
record the dependency chain here so that the headline theorem has a proof at
the point where it is stated.  The tolerant-bisection realizer proved in
\cref{co:thm:bisection-correct} uniformly transforms a fast name of a sequence
$(y_i)$ in $[0,1]$ into a fast name of $(G(y_i))$.  This is
\cref{co:cor:G-sequential}.  Applying the computable, $1$-Lipschitz clamp first
gives the same uniform transformation for $\InvT$ by
\cref{co:lem:clamp-computable,co:cor:GT-sequential}, proving item~1.

For item~2, \cref{co:lem:box-event} gives
$\Delta_r\le F(2^{-r})$.  The exact inverse-modulus theorem, in the
effective-injectivity form \cref{co:cor:effective-injectivity}, therefore gives
\[
 |u-v|<\Delta_r\quad\Longrightarrow\quad
 |\InvT(u)-\InvT(v)|<2^{-r};
\]
this is \cref{co:thm:closed-dyadic-modulus}.  For $n\ge1$, the least integer
$r(n)\ge1$ with $2^{r(n)}>n$ is computable by bounded search, and the displayed
formula for $d(n)$ is a computable composition of integer operations.  Moreover
$d(n)^{-1}=\Delta_{r(n)}$ and $2^{-r(n)}<1/n$.  Chaining these two strict
inequalities with item~2 proves item~3, as detailed in
\cref{co:thm:reciprocal-modulus}.  Thus $G$ (on its subset domain) and $\InvT$
satisfy both clauses of \cref{co:def:computable-function}.
\end{proof}

The theorem is stronger than a nonconstructive assertion that a continuous
inverse on a compact interval is uniformly continuous.  It gives a concrete
integer modulus, a finite algorithm on names, and a sharp structural identity
for the best possible modulus.

\subsection{Four proof layers}

The argument is organized so that every hidden effectiveness issue is exposed.

\begin{enumerate}
\item \textbf{Forward oracle.}  A centered finite convolution has an exact
Thue--Morse truncated-power formula.  At rational inputs it is rational, and
its error from $F$ is at most $2^{-N}$.
\item \textbf{Effective injectivity.}  Symmetric unimodality shows that an
interval of geometric length $h$ has probability mass at least $F(h)$;
strict radial unimodality shows that equality occurs only for a degenerate
interval or an endpoint interval.
A direct event gives the computable lower bound $F(2^{-r})\ge\Delta_r$.
\item \textbf{Uniform continuity.}  Effective injectivity of $F$ becomes an
effective modulus for $G$.  In fact the ordinary modulus of $G$ is exactly
$G$ itself.
\item \textbf{Name transformation.}  Tolerant bisection combines the forward
oracle and the inverse modulus.  A midpoint is used for a left/right update
only when the sign is certified; otherwise it is already a certified output.
\end{enumerate}

The fourth layer matters.  Naive bisection that waits until it decides whether
$F(c)<y$, $F(c)=y$, or $F(c)>y$ is not an algorithm on arbitrary computable
reals: equality of computable reals is undecidable in general.  The tolerance
branch is what removes that obstruction.

% Source fragment: Inverse computability; prefix co.
\section{Computable-analysis definitions}

\subsection{Fast names and computable sequences}

Following standard computable analysis \cite{co:G1957,co:Weihrauch2000,co:PourElRichards1989}, a convenient representation of a real number $x$ is a fast dyadic Cauchy
name: an algorithm returns a dyadic rational $q_p$ such that
\begin{equation}\label{co:eq:fast-name}
 \AbsoluteValue{x-q_p}\le2^{-p}.
\end{equation}
A real sequence $(x_i)_{i\in\NaturalNumbers}$ is computable when one algorithm returns
$q_{i,p}$ satisfying \eqref{co:eq:fast-name} uniformly in both $i$ and $p$.
The repository uses an equivalent signed-natural code for the dyadic
numerator, avoiding an opaque integer encoding.

\begin{definition}[Sequential computability]\label{co:def:sequential}
A function $f$ is \emph{sequentially computable} if there is one effective
realizer transformer with the following property.  Given as a subroutine any
algorithm $Q(i,p)$ returning a dyadic rational with
$|Q(i,p)-x_i|\le2^{-p}$ for a sequence $(x_i)$ in the domain, the transformer
produces an algorithm $R(i,p)$ returning a dyadic rational with
$|R(i,p)-f(x_i)|\le2^{-p}$.  The same transformer works for every input name,
uniformly in the sequence index $i$ and requested precision $p$.
\end{definition}

Thus this is a realizability assertion, not merely the extensional statement
that each computable input sequence happens to have a computable image.  The
construction in \cref{co:sec:sequential} supplies the single uniform
transformer explicitly.

\subsection{The reciprocal uniform-continuity convention}

\begin{definition}[Effective uniform continuity]\label{co:def:euc}
A function $f:\RealNumbers\to\RealNumbers$ is \emph{effectively uniformly continuous} if there
is a recursive $d:\NaturalNumbers\to\NaturalNumbers$ with $d(n)>0$ for $n>0$ such that
\begin{equation}\label{co:eq:euc-def}
 \AbsoluteValue{x-y}<\frac1{d(n)}
 \quad\Longrightarrow\quad
 \AbsoluteValue{f(x)-f(y)}<\frac1n
 \qquad(n>0).
\end{equation}
\end{definition}

This is the reciprocal convention used in the cited Wikipedia article
\cite{co:WikipediaCRF}.  The repository implements it in
\path{FabiusComputability.lean}.  It is equivalent to the more common binary
modulus convention
\begin{equation}\label{co:eq:binary-modulus-convention}
 \AbsoluteValue{x-y}<2^{-m(p)}
 \quad\Longrightarrow\quad
 \AbsoluteValue{f(x)-f(y)}<2^{-p},
\end{equation}
provided $m$ is recursive.  We give both forms because the binary version
reveals the endpoint conditioning much more clearly.

\begin{definition}[Computable real function]\label{co:def:computable-function}
In the Grzegorczyk formulation used here, a real function is computable if it
is sequentially computable and effectively uniformly continuous.
\end{definition}

For $G$ defined only on $[0,1]$, the same definitions are restricted to
sequences and pairs of points in that computable compact domain.  The totalized
$\InvT$ lets one use the literal type $\RealNumbers\to\RealNumbers$ without changing the
mathematics.

\section{Probability, Fourier structure, and the forward function}
\label{co:sec:forward}

\subsection{The random-series model}

The classical probability model may be written as follows\ \cite{is:p1:Fabius1966,is:p1:Arias1982,is:p1:ProveItPrimary}.  Recall that $V_1,V_2,\ldots$ are independent uniform random variables on $[0,1]$ and
put
\begin{equation}\label{co:eq:random-series}
 X:=\sum_{k=1}^{\infty}2^{-k}V_k.
\end{equation}
The series converges absolutely, $0\le X\le1$, and
\begin{equation}\label{co:eq:F-cdf}
 F(x)=\Probability[X\le x].
\end{equation}
Reflection $V_k\mapsto1-V_k$ sends $X$ to $1-X$ and proves the symmetry in
\eqref{co:eq:F-basic}.

The centered random variable
\begin{equation}\label{co:eq:centered-random}
 Y:=2X-1=\sum_{k=1}^{\infty}2^{-k}(2V_k-1)
\end{equation}
has density $\RvachevUp$.  With $\SincRad z=\sin z/z$,
\begin{equation}\label{co:eq:sinc-product}
 \Expectation\EulerE^{\ImaginaryUnit tY}
 =\prod_{k=1}^{\infty}\SincRad\!\left(\frac{t}{2^k}\right).
\end{equation}
Thus the probability model, the Rvachev bump, and the infinite sinc product
are three representations of the same object.  The computability proof uses
the probability representation in physical space; the Fourier and polynomial
alternatives are developed in the shared platform and the deconvolution and
self-sampling parts of this volume.

\subsection{Finite sums and their exact CDFs}

For $N\ge1$, define
\begin{equation}\label{co:eq:finite-random-sum}
 X_N:=\sum_{k=1}^{N}2^{-k}V_k,
 \qquad
 R_N:=X-X_N.
\end{equation}
Then
\begin{equation}\label{co:eq:tail-bound}
 0\le R_N\le2^{-N}.
\end{equation}
Let $F_N(x)=\Probability[X_N\le x]$.  The general inclusion--exclusion formula for a
sum of independent, nonidentically scaled uniforms becomes especially rigid
for dyadic scales.

\begin{proposition}[Finite Thue--Morse CDF]\label{co:prop:finite-cdf}
For $N\ge1$ and $x\in\RealNumbers$,
\begin{equation}\label{co:eq:finite-cdf-formula}
 F_N(x)=
 \frac{2^{N(N+1)/2}}{N!}
 \sum_{j=0}^{2^N-1}(-1)^{\BinaryDigitSum(j)}
 \PositivePart{x-\frac{j}{2^N}}^{N},
\end{equation}
where $\BinaryDigitSum(j)$ is the sum of the binary digits of $j$.
\end{proposition}

\begin{proof}
For positive lengths $a_1,\ldots,a_N$, the CDF of
$\sum a_kV_k$ is
\[
 \frac{1}{N!\prod_{k=1}^N a_k}
 \sum_{A\subseteq\{1,\ldots,N\}}(-1)^{|A|}
 \PositivePart{x-\sum_{k\in A}a_k}^{N}.
\]
This follows by integrating the box indicator over the half-space
$\sum a_ku_k\le x$, or inductively by finite differences of truncated powers.
Here $a_k=2^{-k}$ and
$\prod a_k=2^{-N(N+1)/2}$.  The map
\[
 A\longmapsto j(A):=\sum_{k\in A}2^{N-k}
\]
is a bijection from subsets to $\{0,\ldots,2^N-1\}$, with
$\sum_{k\in A}2^{-k}=j(A)/2^N$ and
$|A|\equiv \BinaryDigitSum(j(A))\pmod2$.  Substitution gives
\eqref{co:eq:finite-cdf-formula}.
\end{proof}

The alternating signs are exactly the Thue--Morse signs.  No limiting
argument is hidden in \eqref{co:eq:finite-cdf-formula}: at rational $x$ it is an
exact rational obtained by finitely many integer operations.

\subsection{Centered finite splines and a uniform certificate}

The deterministic midpoint of the omitted tail is $2^{-N-1}$.  Define
\begin{equation}\label{co:eq:centered-truncation}
 Q_N:=X_N+2^{-N-1},
 \qquad
 S_N(x):=\Probability[Q_N\le x]=F_N(x-2^{-N-1}).
\end{equation}
Then \cref{co:prop:finite-cdf} gives
\begin{equation}\label{co:eq:centered-spline-formula}
 S_N(x)=
 \frac{2^{N(N+1)/2}}{N!}
 \sum_{j=0}^{2^N-1}(-1)^{\BinaryDigitSum(j)}
 \PositivePart{x-\frac{2j+1}{2^{N+1}}}^{N}.
\end{equation}
This is the centered uniform spline used by the repository's primitive-
recursive evaluator.

\begin{lemma}[Uniform density bound]\label{co:lem:density-bound}
The finite CDFs $F_N$ and $S_N$ are globally $2$-Lipschitz.
\end{lemma}

\begin{proof}
The density of $2^{-1}U_1$ is $2\IndicatorOf{[0,1/2]}$.  Convolution with any
probability measure cannot increase the $L^\infty$ norm of a density.  Hence
the density of $X_N$ is bounded by $2$.  Translating by $2^{-N-1}$ does not
change this bound.  Integrating the density over an interval proves the
Lipschitz estimate.
\end{proof}

\begin{theorem}[Certified centered approximation]\label{co:thm:centered-error}
For every $N\ge1$ and every real $x$,
\begin{equation}\label{co:eq:centered-error}
 \AbsoluteValue{S_N(x)-F(x)}\le2^{-N}.
\end{equation}
\end{theorem}

\begin{proof}
Write
\[
 X=Q_N+E_N,
 \qquad
 E_N=R_N-2^{-N-1}.
\]
By \eqref{co:eq:tail-bound}, $|E_N|\le2^{-N-1}$.  Therefore
\[
 S_N(x-2^{-N-1})\le F(x)\le S_N(x+2^{-N-1}).
\]
By \cref{co:lem:density-bound}, moving the argument by $2^{-N-1}$ changes $S_N$
by at most $2\cdot2^{-N-1}=2^{-N}$.  This gives both sides of
\eqref{co:eq:centered-error}.
\end{proof}

\begin{corollary}[Self-contained computability of $F$]\label{co:cor:F-computable}
There is an explicit uniform algorithm which, given a rational $x$ and a
precision $p$, returns a rational $q$ with
$|q-F(x)|\le2^{-p}$.  One may take $q=S_N(x)$ with $N\ge p$.
Consequently $F$ is sequentially computable; together with its global
$2$-Lipschitz bound, it is a computable real function.
\end{corollary}

\begin{proof}
Formula \eqref{co:eq:centered-spline-formula} is a finite rational computation,
and \cref{co:thm:centered-error} supplies its error certificate.  At an arbitrary
computable input, first replace the input by a sufficiently accurate dyadic;
the $2$-Lipschitz bound controls this replacement.  The construction is
uniform in a sequence index.  Effective uniform continuity follows from
$|F(x)-F(y)|\le2|x-y|$; for example the reciprocal modulus $d(n)=2n$ works.
\end{proof}

\begin{remark}[Efficiency versus computability]
The direct sum has $2^N$ terms and is not intended as an efficient evaluator.
The repository's exact dyadic recurrence, highest-bit reduction, and
primitive-recursive natural-number implementation are substantially better.
For the theorem here, only termination and a proved error bound are required.
\end{remark}

\section{Effective positivity, strict increase, and the density shape}

\subsection{A closed positive lower bound at dyadic points}

The existence of an inverse requires strict increase; effective inversion
requires a quantitative version.  Both begin with an elementary positive
probability event.

\begin{lemma}[Dyadic endpoint box event]\label{co:lem:box-event}
For every $r\ge0$,
\begin{equation}\label{co:eq:Delta-def}
 F(2^{-r})\ge
 \Delta_r:=2^{-r(r+1)/2}(r+1)^{-(r+1)}.
\end{equation}
In particular, $F(x)>0$ for every $x>0$.
\end{lemma}

\begin{proof}
Put $m=r+1$.  For $1\le k\le m$, impose the independent restrictions
\begin{equation}\label{co:eq:box-restrictions}
 V_k\le a_k:=\frac{2^{k-r-1}}{m}.
\end{equation}
The largest $a_k$ is $a_m=1/m\le1$, so these are valid events.  On their
intersection,
\[
 2^{-k}V_k\le\frac{2^{-r-1}}m
 \quad(1\le k\le m).
\]
After summation the first $m$ terms contribute at most $2^{-r-1}$.  The
unrestricted tail contributes at most
\[
 \sum_{k=m+1}^{\infty}2^{-k}=2^{-m}=2^{-r-1}.
\]
Thus the event in \eqref{co:eq:box-restrictions} implies $X\le2^{-r}$.  Its
probability is
\begin{align*}
 \prod_{k=1}^{m}a_k
 &=m^{-m}2^{\sum_{k=1}^m(k-r-1)}\\
 &=m^{-m}2^{m(m+1)/2-m(r+1)}\\
 &=2^{-r(r+1)/2}(r+1)^{-(r+1)}=\Delta_r.
\end{align*}
This proves \eqref{co:eq:Delta-def}.  Given $x>0$, choose $r$ with
$2^{-r}\le x$ and use monotonicity of a CDF.
\end{proof}

The bound is deliberately elementary: it uses only independence, the support
of the tail, and integer arithmetic.  It is not numerically sharp, but
\cref{co:sec:binary-conditioning} shows that its logarithm has the correct two
largest orders for computability purposes.

The numerical conclusion of \cref{co:lem:box-event} is now Lean-formalized, but
by a different and stronger route.  Here is the complete recurrence argument.
Let $a_0=1$ and use the repository's exact inverse-dyadic identities
\[
 a_r=
 \frac{\displaystyle\sum_{0\le k<r}a_k/(r-k+1)!}{2^r-1}
 \quad(r\ge1),
 \qquad
 F(2^{-r})=2^{-\binom r2}a_r.
\]
Every $a_k$ is positive.  Retaining only the $k=0$ summand therefore gives
\[
 a_r\ge\frac{1}{(r+1)!(2^r-1)},
 \qquad
 F(2^{-r})\ge
 \frac{1}{2^{\binom r2}(r+1)!(2^r-1)}
 \quad(r\ge1).
\]
Since $2^r-1\le2^r$ and
$\binom r2+r=\binom{r+1}{2}$, taking the larger denominator yields the
first inequality below; the exceptional index $r=0$ is simply $F(1)=1$.
Thus for every $r\ge0$,
\begin{equation}\label{co:eq:factorial-endpoint-lower}
 F(2^{-r})\ge
 \frac{1}{2^{\binom{r+1}{2}}(r+1)!}
 \ge
 \frac{1}{2^{\binom{r+1}{2}}(r+1)^{r+1}}
 =\Delta_r.
\end{equation}
Indeed, $(r+1)!\le(r+1)^{r+1}$ proves the second inequality after
multiplication by the positive power of two and reversal under reciprocation.

For the executable arithmetic, write
\[
 d_{\mathrm{fac}}(r):=2^{\binom{r+1}{2}}(r+1)!,
 \qquad
 d_{\Delta}(r):=2^{\binom{r+1}{2}}(r+1)^{r+1}=\Delta_r^{-1}.
\]
Lean defines the first denominator by primitive recursion,
\[
 d_{\mathrm{fac}}(0)=1,
 \qquad
 d_{\mathrm{fac}}(r+1)
 =d_{\mathrm{fac}}(r)\,2^{r+1}(r+2).
\]
Induction, using
$\binom{r+2}{2}=\binom{r+1}{2}+r+1$, proves the displayed closed form.
The recursive clause uses only successor, multiplication, and exponentiation,
so $d_{\mathrm{fac}}$ is primitive recursive.  Likewise
$\binom{r+1}{2}=r(r+1)/2$, and primitive-recursive closure under successor,
multiplication, natural-number division, and exponentiation proves that
$d_{\Delta}$ is primitive recursive.  Finally
$d_{\mathrm{fac}}(r)\le d_{\Delta}(r)$ is exactly the factorial--power
inequality above, and positivity reverses this comparison between their
reciprocals.

The declarations
\path{Fabius.fabiusReal_inverse_two_pow_one_term_lower_bound},
\path{Fabius.inv_inverseFabiusFactorialDenominator_le_fabiusReal}, and
\path{Fabius.inv_inverseFabiusDeltaDenominator_le_fabiusReal} certify this
chain.  They do not formalize the event in
\eqref{co:eq:box-restrictions}, its containment, or its product probability;
that probabilistic proof remains report-level.  The argument in this paragraph
uses only the already formalized dyadic recurrence and therefore preserves the
distinction between the two proof routes.

\subsection{Strict monotonicity from the random series}

\begin{proposition}[Strict increase]\label{co:prop:strict-increase}
For $0\le a<b\le1$,
\begin{equation}\label{co:eq:strict-increase}
 F(a)<F(b).
\end{equation}
Hence $F:[0,1]\to[0,1]$ is a homeomorphism and $G=F^{-1}$ exists.
\end{proposition}

\begin{proof}
The density of $X_N$ is strictly positive on the interior of its support
$(0,1-2^{-N})$.  This follows inductively: a convolution of two densities
that are positive on the interiors of compact intervals is positive on the
interior of their Minkowski-sum interval, because the convolution integral
contains a nonempty open region on which both factors are positive.

Choose $N$ so large that $2^{-N}<b-a$.  Then the open interval
$(a,b-2^{-N})$ is nonempty and lies below the right endpoint of the support of
$X_N$.  Choose a nonempty open subinterval $J$ of it.  The preceding positivity
gives $\Probability[X_N\in J]>0$.  For $X_N\in J$ and every possible
$0\le R_N\le2^{-N}$, one has $a<X_N+R_N<b$.  Consequently
\[
 F(b)-F(a)=\Probability[a<X\le b]\ge\Probability[X_N\in J]>0.
\]
Continuity and the endpoint values $F(0)=0$, $F(1)=1$ then give the
homeomorphism statement.
\end{proof}

\subsection{Symmetric strict unimodality}

Let
\begin{equation}\label{co:eq:density-p}
 p(x):=F'(x),\qquad 0\le x\le1.
\end{equation}
The repository proves that $F$ is $C^\infty$, so $p$ is continuous.  The
functional and reflection equations completely determine its qualitative
shape.

\begin{proposition}[Shape of the Fabius density]\label{co:prop:density-shape}
The density $p$ satisfies
\begin{equation}\label{co:eq:density-symmetry}
 p(x)=p(1-x),
\end{equation}
is strictly increasing on $[0,1/2]$, and is strictly decreasing on
$[1/2,1]$.  Equivalently, $p(x)$ is a strictly decreasing function of
$|x-1/2|$.
\end{proposition}

\begin{proof}
On $[0,1/2]$, \eqref{co:eq:F-basic} gives
\[
 p(x)=2F(2x).
\]
By \cref{co:prop:strict-increase}, this is strictly increasing.  Differentiating
$F(x)+F(1-x)=1$ gives $p(x)=p(1-x)$, which reflects the first-half behavior
to the second half.
\end{proof}

This simple observation is the geometric heart of the inverse modulus.  It is
also a point where Rvachev's function enters directly: by
\eqref{ii:eq:up-F},
\begin{equation}\label{co:eq:p-up}
 p(x)=2\RvachevUp(2x-1).
\end{equation}
Thus \cref{co:prop:density-shape} is the standard single-peaked shape of the
compactly supported atomic bump.

\section{The least-mass interval and the exact inverse modulus}
\label{co:sec:exact-modulus}

\subsection{Intervals of a prescribed geometric length}

For $0\le h\le1$ and $0\le x\le1-h$, define the fixed-length CDF increment
\begin{equation}\label{co:eq:interval-mass}
 M_h(x):=F(x+h)-F(x).
\end{equation}
If $\mu$ is the Fabius probability law with CDF $F$, then the canonical
measure-theoretic interpretation is
\[
 M_h(x)=\mu\bigl((x,x+h]\bigr).
\]
The law is atomless---indeed it has the continuous density $p=F'$---so the
value is unchanged if either or both endpoints are included.  In particular,
\[
 \mu\bigl((x,x+h]\bigr)
 =\mu\bigl([x,x+h]\bigr)
 =\mu\bigl((x,x+h)\bigr)
 =\mu\bigl([x,x+h)\bigr).
\]
We therefore continue to call $M_h(x)$ an interval mass, while
\eqref{co:eq:interval-mass} fixes the canonical CDF convention.

\begin{theorem}[Strict interval-mass shape and exact minimizers]
\label{co:thm:least-mass}
For every $0\le h\le1$,
\begin{equation}\label{co:eq:least-mass}
 \min_{0\le x\le1-h}\bigl(F(x+h)-F(x)\bigr)=F(h).
\end{equation}
If $0<h\le1$ and $m=(1-h)/2$, then the globally extended increment $M_h$ is
strictly increasing on $[-h,m]$, symmetric about $m$, and strictly decreasing
on $[m,1]$; outside $[-h,1]$ it is identically zero.  On the admissible-start
interval $0\le x\le1-h$ one has the boundary-safe equality law
\begin{equation}\label{co:eq:least-mass-equality}
 M_h(x)=F(h)
 \quad\Longleftrightarrow\quad
 h=0\ \text{or}\ x=0\ \text{or}\ x=1-h.
\end{equation}
Equivalently, whenever $0\le a\le b\le1$,
\begin{equation}\label{co:eq:least-mass-equality-ab}
 F(b-a)=F(b)-F(a)
 \quad\Longleftrightarrow\quad
 a=b\ \text{or}\ a=0\ \text{or}\ b=1.
\end{equation}
Thus every genuinely interior interval $0<a<b<1$ has strictly more Fabius
mass than an endpoint interval of the same length.  At $h=1$ the two endpoint
formulas coincide at the single admissible start $x=0$.
\end{theorem}

\begin{proof}
The case $h=0$ gives $M_0\equiv0$, so suppose $0<h\le1$ and put
$m=(1-h)/2$.  Reflection gives
\begin{align*}
 M_h(1-h-x)
 &=F(1-x)-F(1-h-x)\\
 &=(1-F(x))-(1-F(x+h))=M_h(x),
\end{align*}
so $M_h$ is symmetric about $m$.

Consider $-h<x<m$.  Then $0<x+h<1$.  If $x\le0$, the constant left tail gives
$p(x)=0$, whereas strict positivity on the open support gives $p(x+h)>0$.
If $x>0$, compute
\begin{equation}\label{co:eq:distance-comparison}
 \left(x+h-\frac12\right)^2-
 \left(x-\frac12\right)^2
 =h(2x+h-1)<0.
\end{equation}
Thus $x+h$ is strictly closer to $1/2$ than $x$ is.  The strict radial
unimodality in \cref{co:prop:density-shape} again gives $p(x+h)>p(x)$.  Hence
\begin{equation}\label{co:eq:M-derivative}
 M_h'(x)=p(x+h)-p(x)>0
 \qquad(-h<x<m).
\end{equation}
The mean-value theorem makes $M_h$ strictly increasing on the closed interval
$[-h,m]$.  Reflection makes it strictly decreasing on $[m,1]$.  The clamped
tails of $F$ give $M_h=0$ to the left of $-h$ and to the right of $1$; this is
why the global half-line shape is only non-strict, whereas the two finite
branches above are maximally strict.

Finally,
\[
 M_h(0)=F(h)-F(0)=F(h)
 \quad\text{and}\quad
 M_h(1-h)=F(1)-F(1-h)=F(h).
\]
If $0<x<1-h$, compare $x$ strictly with $0$ on the left branch when $x\le m$,
and compare it strictly with $1-h$ on the right branch when $m<x$.  In either
case $M_h(x)>F(h)$.  This proves the minimum and the complete equality locus
\eqref{co:eq:least-mass-equality}.  Substituting $h=b-a$ and $x=a$ gives
\eqref{co:eq:least-mass-equality-ab}; its diagonal alternative $a=b$ is exactly
the degenerate case $h=0$.
\end{proof}

\begin{corollary}[Global increment shape and constrained superadditivity]
\label{co:cor:global-increment-shape}
For every $h\ge0$, the globally extended increment $M_h$ is nondecreasing on
the half-line
\[
 x\le\frac{1-h}{2}
\]
and nonincreasing on the complementary half-line.  Consequently, whenever
$a,b\ge0$ and $a+b\le1$,
\begin{equation}\label{co:eq:fabius-superadditive}
 F(a)+F(b)\le F(a+b).
\end{equation}
Its complete equality locus is
\begin{equation}\label{co:eq:fabius-superadditive-equality}
 F(a)+F(b)=F(a+b)
 \quad\Longleftrightarrow\quad
 a=0\ \text{or}\ b=0\ \text{or}\ a+b=1.
\end{equation}
\end{corollary}

\begin{proof}
The derivative argument above never used $x\ge0$: from
$x\le(1-h)/2$ and $h\ge0$, either $x+h\le1/2$ and the monotonicity of $p$ on
the left half gives $p(x)\le p(x+h)$, or $x+h>1/2$ and reflection replaces
$p(x+h)$ by $p(1-x-h)$, where
\[
 x\le1-x-h\le\frac12.
\]
Thus $M_h'\ge0$ on the entire left half-line.  Reflection
$M_h(1-h-x)=M_h(x)$ gives the right-half statement.  For
\eqref{co:eq:fabius-superadditive}, apply \cref{co:thm:least-mass} to the interval
from $a$ to $a+b$: its length is $b$, so
\[
 F(b)\le F(a+b)-F(a).
\]
 Rearranging proves the inequality.  Equality is exactly the equality case of
 \cref{co:thm:least-mass} for the interval $[a,a+b]$: its three alternatives are
 $a=a+b$, $a=0$, and $a+b=1$, and the first is equivalent to $b=0$.
\end{proof}

\begin{remark}[A general principle]\label{co:rem:general-principle}
The lower bound used only that a probability density on $[0,1]$ is symmetric
and unimodal.  Hence the minimum-value part of \cref{co:thm:least-mass} holds for
every continuous symmetric unimodal distribution.  The endpoint-only equality
locus uses the genuinely stronger hypothesis that the density decreases
strictly with distance from the midpoint; plateaus can create additional
minimizers.  The Fabius system is special not because the geometric lemma is
exotic, but because $F(h)$ has unusually rich exact arithmetic and endpoint
asymptotics.
\end{remark}

\subsection{The inverse's modulus is the inverse itself}

For a bounded function $H$ define its ordinary modulus of continuity by
\begin{equation}\label{co:eq:ordinary-modulus}
 \omega_H(\eta):=
 \sup\SetOf{\AbsoluteValue{H(u)-H(v)}:\AbsoluteValue{u-v}\le\eta}.
\end{equation}
For $G$ the variables are restricted to $[0,1]$; for $\InvT$ they range over
$\RealNumbers$.

\begin{theorem}[Exact modulus of the inverse]\label{co:thm:exact-inverse-modulus}
For every $\eta\ge0$, with the inputs in the definition of $\omega_G$
restricted to $[0,1]$,
\begin{equation}\label{co:eq:exact-inverse-modulus}
 \omega_G(\eta)=G(\min(\eta,1)).
\end{equation}
In particular this is $G(\eta)$ for $0\le\eta\le1$ and is $1$ thereafter.
For the totalized inverse and every $\eta\ge0$,
\begin{equation}\label{co:eq:exact-total-modulus}
 \omega_{\InvT}(\eta)=G(\min(\eta,1)).
\end{equation}
More strongly, for all real $u,v$,
\begin{equation}\label{co:eq:absolute-inverse-gap-bound}
 \AbsoluteValue{\InvT(u)-\InvT(v)}
 \le \InvT(\AbsoluteValue{u-v})
 =G\bigl(\min(\AbsoluteValue{u-v},1)\bigr).
\end{equation}
The associated ordered-gap and subadditivity laws are
\begin{equation}\label{co:eq:inverse-gap-bound}
 \InvT(v)-\InvT(u)\le \InvT(v-u)
 \qquad(u,v\in\RealNumbers)
\end{equation}
and
\begin{equation}\label{co:eq:inverse-subadditive}
 \InvT(x+y)\le\InvT(x)+\InvT(y)
 \qquad(x,y\in\RealNumbers).
\end{equation}
On the unit interval the ordered-gap equality cases are exactly
\begin{equation}\label{co:eq:inverse-gap-equality}
 G(v)-G(u)=G(v-u)
 \quad\Longleftrightarrow\quad
 u=v\ \text{or}\ u=0\ \text{or}\ v=1
 \qquad(0\le u\le v\le1).
\end{equation}
Equivalently, without ordering the inputs,
\begin{equation}\label{co:eq:absolute-inverse-gap-equality}
 \AbsoluteValue{G(u)-G(v)}=G(\AbsoluteValue{u-v})
 \quad\Longleftrightarrow\quad
 u=v\ \text{or}\ u\in\{0,1\}\ \text{or}\ v\in\{0,1\}
 \qquad(u,v\in[0,1]).
\end{equation}
These two equality classifications are intentionally unit-domain statements:
the constant tails of $\InvT$ create additional configurations on $\RealNumbers$.
\end{theorem}

\begin{proof}
First let $c,d\in[0,1]$.  Interchanging them if necessary, assume $c\le d$,
and put
\[
 a=G(c),\qquad b=G(d).
\]
Monotonicity of $G$ gives $a\le b$.  The inverse identities and
\cref{co:thm:least-mass} give
\begin{equation}\label{co:eq:mass-gap-proof}
 F(b-a)\le F(b)-F(a)=d-c.
\end{equation}
Both $b-a$ and $d-c$ lie in $[0,1]$.  Applying the increasing inverse and
using $G(F(b-a))=b-a$ yields
\[
 b-a\le G(d-c).
\]
Removing the temporary ordering assumption gives
\begin{equation}\label{co:eq:unit-absolute-gap}
 \AbsoluteValue{G(c)-G(d)}\le G(\AbsoluteValue{c-d})
 \qquad(c,d\in[0,1]).
\end{equation}
The same argument records equality.  Under $c\le d$, equality in the inverse
gap is equivalent, after applying the strictly increasing $F$, to equality in
\eqref{co:eq:mass-gap-proof}.  By \eqref{co:eq:least-mass-equality-ab}, this says
\[
 a=b,\qquad a=0,\qquad\text{or}\qquad b=1.
\]
Because $a=G(c)$ and $b=G(d)$ and $F,G$ are inverse order isomorphisms, these
alternatives are respectively $c=d$, $c=0$, and $d=1$.  This proves
\eqref{co:eq:inverse-gap-equality}.  Applying it after ordering $c,d$ proves
\eqref{co:eq:absolute-inverse-gap-equality}; reversing the order exchanges the
left and right endpoint alternatives.

Now put $c=\Clamp(u)$ and $d=\Clamp(v)$.  Projection onto $[0,1]$ is
nonexpansive and has diameter one, so
\[
 \AbsoluteValue{c-d}\le\min(\AbsoluteValue{u-v},1).
\]
Combining this with \eqref{co:eq:unit-absolute-gap} proves
\eqref{co:eq:absolute-inverse-gap-bound}.  If $\AbsoluteValue{u-v}\le\eta$, monotonicity
then gives the upper bound $G(\min(\eta,1))$ for both moduli.  Let
$\theta=\min(\eta,1)$.  The pair $(0,\theta)$ belongs to $[0,1]^2$, has input
separation $\theta\le\eta$, and has output separation
\[
 \AbsoluteValue{G(\theta)-G(0)}=G(\theta).
\]
Thus the displayed upper bound is attained, not merely approached, proving
both exact supremum identities.

For \eqref{co:eq:inverse-gap-bound}, if $u\le v$, remove the absolute value from
\eqref{co:eq:absolute-inverse-gap-bound}; then $v-u=\AbsoluteValue{v-u}$.  If $v<u$, the
left side is nonpositive by monotonicity, whereas $\InvT(v-u)\ge0$.
Finally apply \eqref{co:eq:inverse-gap-bound} with $u=x$ and $v=x+y$:
\[
 \InvT(x+y)-\InvT(x)\le\InvT(y),
\]
which rearranges to \eqref{co:eq:inverse-subadditive}.
\end{proof}

\begin{corollary}[Effective-injectivity form]\label{co:cor:effective-injectivity}
For $0\le h\le1$ and $u,v\in\RealNumbers$,
\begin{equation}\label{co:eq:effective-injectivity}
 \AbsoluteValue{u-v}<F(h)
 \quad\Longrightarrow\quad
 \AbsoluteValue{\InvT(u)-\InvT(v)}<h.
\end{equation}
The closed-threshold and contrapositive forms are
\begin{align}
 \AbsoluteValue{u-v}\le F(h)
 &\Longrightarrow \AbsoluteValue{\InvT(u)-\InvT(v)}\le h,
 \label{co:eq:closed-effective-injectivity}\\
 h\le\AbsoluteValue{\InvT(u)-\InvT(v)}
 &\Longrightarrow F(h)\le\AbsoluteValue{u-v}.
 \label{co:eq:contrapositive-effective-injectivity}
\end{align}
The same statements hold for $G$ on $[0,1]$.
\end{corollary}

\begin{proof}
The strict hypothesis gives
$\AbsoluteValue{u-v}<F(h)\le1$, so the saturation in the exact modulus disappears.
Strict monotonicity of $G$ on $[0,1]$ therefore gives
\[
 \AbsoluteValue{\InvT(u)-\InvT(v)}
 \le G(\AbsoluteValue{u-v})<G(F(h))=h.
\]
The same chain with non-strict inequalities proves
\eqref{co:eq:closed-effective-injectivity}; its contrapositive is
\eqref{co:eq:contrapositive-effective-injectivity}.
\end{proof}

\begin{theorem}[Exact strict threshold]\label{co:thm:exact-strict-threshold}
For $0\le h\le1$ and every real $\delta$,
\begin{equation}\label{co:eq:exact-strict-threshold}
 \left[
   \forall u,v\in\RealNumbers,
   \ \AbsoluteValue{u-v}<\delta
   \Longrightarrow
   \AbsoluteValue{\InvT(u)-\InvT(v)}<h
 \right]
 \quad\Longleftrightarrow\quad
 \delta\le F(h).
\end{equation}
Thus $F(h)$ is exactly the largest strict universal input threshold for output
error below $h$.
\end{theorem}

\begin{proof}
If $\delta\le F(h)$, then
$\AbsoluteValue{u-v}<\delta$ implies $\AbsoluteValue{u-v}<F(h)$, and
\cref{co:cor:effective-injectivity} applies.  Conversely, suppose the universal
implication holds but $F(h)<\delta$.  Take
\[
 u=0,\qquad v=F(h).
\]
Their input separation is $F(h)<\delta$, whereas the inverse identities give
\[
 \AbsoluteValue{\InvT(0)-\InvT(F(h))}=\AbsoluteValue{0-h}=h,
\]
contradicting the asserted strict output bound.  Notice that no sign
hypothesis on $\delta$ is needed: for $\delta\le0$ the left implication is
vacuous and the right inequality follows from $F(h)\ge0$.
\end{proof}

\section{Effective uniform continuity}

\subsection{A closed primitive-recursive modulus}

Combining the exact inverse modulus with the box event gives the promised
fully explicit implication.

\begin{theorem}[Closed dyadic modulus]\label{co:thm:closed-dyadic-modulus}
For every $r\ge0$ and all $u,v\in\RealNumbers$,
\begin{equation}\label{co:eq:closed-dyadic-modulus}
 \AbsoluteValue{u-v}<\Delta_r
 \quad\Longrightarrow\quad
 \AbsoluteValue{\InvT(u)-\InvT(v)}<2^{-r},
\end{equation}
where $\Delta_r$ is given by \eqref{co:eq:Delta-def}.
\end{theorem}

\begin{proof}
By \cref{co:lem:box-event}, $\Delta_r\le F(2^{-r})$.  Apply
\cref{co:cor:effective-injectivity} with $h=2^{-r}$.
\end{proof}

The exact statement \eqref{co:eq:closed-dyadic-modulus} is also formalized as
\path{Fabius.abs_fabiusInv_sub_lt_inverse_two_pow_of_lt_deltaDenominator}.
The stronger factorial-denominator version is
\path{Fabius.abs_fabiusInv_sub_lt_inverse_two_pow_of_lt_factorialDenominator}.
The formal API also closes both endpoints:
\[
 \AbsoluteValue{u-v}\le\Delta_r
 \quad\Longrightarrow\quad
 \AbsoluteValue{\InvT(u)-\InvT(v)}\le2^{-r},
\]
with declaration
\path{Fabius.abs_fabiusInv_sub_le_inverse_two_pow_of_le_deltaDenominator};
\path{Fabius.abs_fabiusInv_sub_le_inverse_two_pow_of_le_factorialDenominator}
is its stronger factorial-denominator companion.  Thus the strict and closed
input/output boundaries are separately available rather than recovered by an
unstated limiting argument.  Their common proof is direct.  Let $d$ be either
$d_{\mathrm{fac}}(r)$ or $d_{\Delta}(r)$.  The recurrence argument above gives
$d^{-1}\le F(2^{-r})$.  Hence
\[
 \AbsoluteValue{u-v}<d^{-1}
 \quad\Longrightarrow\quad
 \AbsoluteValue{u-v}<F(2^{-r})
 \quad\Longrightarrow\quad
 \AbsoluteValue{\InvT(u)-\InvT(v)}<2^{-r}
\]
by strict effective injectivity.  Replacing both strict inequalities by weak
ones and using its closed-threshold companion proves the two closed versions.
This proves all four formal dyadic implications without a limiting argument;
only the choice of denominator distinguishes the factorial and Delta forms.

\begin{theorem}[Reciprocal modulus required by the definition]\label{co:thm:reciprocal-modulus}
For $n\ge1$, let $r(n)$ be the least positive integer with $2^{r(n)}>n$ and
put
\begin{equation}\label{co:eq:d-explicit}
 d(n)=2^{r(n)(r(n)+1)/2}\bigl(r(n)+1\bigr)^{r(n)+1},
 \qquad d(0)=1.
\end{equation}
Then $d$ is recursive and satisfies \eqref{co:eq:euc-def} for $\InvT$.
Thus $\InvT$ and the restricted inverse $G$ are effectively uniformly
continuous.
\end{theorem}

\begin{proof}
The function $r(n)$ can be found by bounded integer search, or equivalently by
computing the binary length of $n$; hence it is recursive.  Exponentiation,
multiplication, and composition preserve recursiveness, so $d$ is recursive.
By construction,
\[
 \frac1{d(n)}=\Delta_{r(n)}
 \quad\text{and}\quad
 2^{-r(n)}<\frac1n.
\]
If $|u-v|<1/d(n)$, \cref{co:thm:closed-dyadic-modulus} gives
$|\InvT(u)-\InvT(v)|<2^{-r(n)}<1/n$.
\end{proof}

\begin{remark}[An even simpler, much larger modulus]
Taking $r=n$ gives the one-line primitive-recursive choice
\[
 d_0(n)=2^{n(n+1)/2}(n+1)^{n+1}\qquad(n\ge1),
\]
because $2^{-n}<1/n$.  The logarithmic choice $r(n)$ in
\cref{co:thm:reciprocal-modulus} is far smaller and has the correct leading
complexity scale.

The Lean development uses the stronger one-line witness
\[
 d_{\mathrm{fac}}(n)=2^{\binom{n+1}{2}}(n+1)!,
\]
which is no larger than $d_0(n)$.  The theorem
\path{Fabius.fabiusInv_effectivelyUniformContinuous} proves
\texttt{EffectivelyUniformContinuous} for the totalized inverse as follows.
The recurrence above makes $d_{\mathrm{fac}}$ primitive recursive and positive.
For $n\ge1$, the strict factorial implication gives
\[
 \AbsoluteValue{u-v}<\frac1{d_{\mathrm{fac}}(n)}
 \quad\Longrightarrow\quad
 \AbsoluteValue{\InvT(u)-\InvT(v)}<2^{-n}.
\]
The elementary inequality $n<2^n$ between positive integers reverses under
reciprocation, so $2^{-n}<1/n$.  Chaining these inequalities is precisely the
reciprocal convention in \cref{co:def:euc}; $d_{\mathrm{fac}}(0)=1$ supplies the
harmless zeroth value.  Thus effective uniform continuity itself is
formalized.  The downstream
\path{FabiusInverseLogarithmicModulus.lean} now also formalizes the smaller
choice $r(n)$ in \cref{co:thm:reciprocal-modulus}.  Its declaration
\path{Fabius.inverseFabiusLogarithmicOrder_isLeast} identifies the least order,
\path{Fabius.inverseFabiusLogarithmicOrder_primrec} proves primitive
recursiveness, and
\path{Fabius.fabiusInv_effectivelyUniformContinuous_logarithmicDelta} packages
the report-exact denominator $d(n)$.  The same module supplies the stronger
logarithmic factorial witness.
\end{remark}

\subsection{The exact-dyadic repository modulus}

Put $h_n:=2^{-r(n)}$.  The corpus proves that every $F(h_n)$ is an exactly
computable positive rational.  Therefore one can replace the elementary lower
bound $\Delta_{r(n)}$ by the exact endpoint mass while retaining an exact
dyadic output proxy.

\begin{proposition}[Repository-native exact-dyadic modulus]\label{co:prop:exact-dyadic-modulus}
Let
\begin{equation}\label{co:eq:d-star}
 d_*(n):=
 \Ceiling{\frac{1}{F(h_n)}}
 \qquad(n\ge1),
 \qquad h_n=2^{-r(n)},
 \quad d_*(0):=1.
\end{equation}
Then $d_*$ is recursive and
\begin{equation}\label{co:eq:d-star-strong}
 \AbsoluteValue{u-v}<\frac1{d_*(n)}
 \quad\Longrightarrow\quad
 \AbsoluteValue{\InvT(u)-\InvT(v)}<h_n<\frac1n.
\end{equation}
Moreover, $d_*(n)$ is the least positive integer $d$ satisfying
\[
 \frac1d\le F(h_n).
\]
Equivalently, it is the least reciprocal denominator certified by the exact
inverse-modulus law for the fixed dyadic output target $h_n$.  It is not
claimed to be the least integer modulus for the weaker target $1/n$.
\end{proposition}

\begin{proof}
The exact dyadic evaluator computes the positive rational $F(h_n)$, and
rational reciprocal and ceiling are computable integer operations.  By the
defining property of the ceiling,
\[
 \frac1{d_*(n)}\le F(h_n).
\]
The strict input hypothesis and \cref{co:cor:effective-injectivity} therefore
give output separation strictly below $h_n$, while the definition of $r(n)$
gives $h_n<1/n$.

For the qualified minimality statement, let $1\le d<d_*(n)$.  Then
\[
 d<\frac1{F(h_n)},
 \qquad\text{hence}\qquad
 \frac1d>F(h_n).
\]
The endpoint pair $u=0$, $v=F(h_n)$ has input separation strictly below
$1/d$ and output separation exactly $h_n$.  Thus $d$ cannot certify the
stricter conclusion $\AbsoluteValue{\InvT(u)-\InvT(v)}<h_n$.
\end{proof}

\begin{remark}[The genuine integer optimum for the $1/n$ target]
The exact inverse-modulus law shows that a positive integer $d$ has the
universal property
\[
 \AbsoluteValue{u-v}<\frac1d
 \quad\Longrightarrow\quad
 \AbsoluteValue{\InvT(u)-\InvT(v)}<\frac1n
\]
if and only if
\[
 \frac1d\le F(1/n).
\]
Consequently the mathematically least such denominator is
\[
 d_{\mathrm{opt}}(n)=\Ceiling{\frac1{F(1/n)}}.
\]
The present report does not prove that $n\mapsto d_{\mathrm{opt}}(n)$ is
recursive: computability of the real number $F(1/n)$ alone does not make the
discontinuous ceiling operation effective at an unknown integer boundary.
The dyadic quantity $d_*$ avoids that issue because $F(h_n)$ is an exact
rational.

The distinction is already visible at $n=1$.  Here $r(1)=1$,
$F(1/2)=1/2$, and hence $d_*(1)=2$, whereas $d=1$ already satisfies the
required implication with output tolerance $1$.
\end{remark}

The exact-ceiling construction $d_*$, including its recursion and qualified
denominator minimality, has not yet been translated into Lean.  The current
logarithmic effective-continuity layer deliberately avoids the ceiling by
using the explicit Delta or stronger factorial denominator at the formalized
least order $r(n)$.

\subsection{Binary conditioning is quadratic}
\label{co:sec:binary-conditioning}

The exact modulus permits a sharp, representation-independent measurement of
how many input bits are intrinsically needed near the flat endpoint.

\begin{definition}[Optimal binary modulus index]\label{co:def:mu}
For $p\ge1$, define
\begin{equation}\label{co:eq:mu-def}
 \mu(p):=\min\SetOf{m\in\NaturalNumbers:2^{-m}\le F(2^{-p})}
 =\Ceiling{\log_2\frac1{F(2^{-p})}}.
\end{equation}
Then
\begin{equation}\label{co:eq:mu-property}
 \AbsoluteValue{u-v}<2^{-\mu(p)}
 \quad\Longrightarrow\quad
 \AbsoluteValue{\InvT(u)-\InvT(v)}<2^{-p},
\end{equation}
and no smaller $m$ has this universal property.
\end{definition}

The lower event bound already gives an upper bound for $\mu(p)$.  A matching
lower bound follows from endpoint flatness.

\begin{lemma}[Elementary flatness upper bound]\label{co:lem:flatness-upper}
For every integer $q\ge1$ and $0\le x\le2^{-q}$,
\begin{equation}\label{co:eq:flatness-upper}
 F(x)\le\frac{2^{q(q+1)/2}}{q!}\,x^q.
\end{equation}
In particular,
\begin{equation}\label{co:eq:dyadic-flatness-upper}
 F(2^{-p})\le\frac{2^{-p(p-1)/2}}{p!}.
\end{equation}
\end{lemma}

\begin{proof}
Repeated differentiation of $F'(x)=2F(2x)$ gives
\begin{equation}\label{co:eq:iterated-derivative}
 F^{(q)}(x)=2^{q(q+1)/2}F(2^qx)
 \qquad(0\le x\le2^{-q}).
\end{equation}
All derivatives of $F$ vanish at $0$, and $0\le F\le1$.  Taylor's formula in
integral form therefore gives
\begin{align*}
 F(x)
 &=\frac1{(q-1)!}\int_0^x(x-t)^{q-1}F^{(q)}(t)\Differential t\\
 &\le\frac{2^{q(q+1)/2}}{(q-1)!}
      \int_0^x(x-t)^{q-1}\Differential t
 =\frac{2^{q(q+1)/2}}{q!}x^q.
\end{align*}
Set $q=p$ and $x=2^{-p}$ to obtain
\eqref{co:eq:dyadic-flatness-upper}.
\end{proof}

\begin{theorem}[Asymptotically sharp input-bit requirement]\label{co:thm:mu-asymptotic}
For every $p\ge1$,
\begin{align}
 \Ceiling{\frac{p(p-1)}2+\log_2(p!)}
 \;&\le\;\mu(p)\label{co:eq:mu-lower}\\
 &\le\;
 \Ceiling{\frac{p(p+1)}2+(p+1)\log_2(p+1)}.
 \label{co:eq:mu-upper}
\end{align}
Consequently
\begin{equation}\label{co:eq:mu-asymptotic}
 \boxed{\displaystyle
 \mu(p)=\frac{p^2}{2}+p\log_2p+\BigO(p).}
\end{equation}
\end{theorem}

\begin{proof}
The flatness bound \eqref{co:eq:dyadic-flatness-upper} yields
\[
 \log_2\frac1{F(2^{-p})}
 \ge\frac{p(p-1)}2+\log_2(p!),
\]
while \cref{co:lem:box-event} yields
\[
 \log_2\frac1{F(2^{-p})}
 \le\frac{p(p+1)}2+(p+1)\log_2(p+1).
\]
Take ceilings and use \eqref{co:eq:mu-def}.  Stirling's formula
$\log_2(p!)=p\log_2p-(\log_2\EulerE)p+\BigO(\log p)$ makes both bounds equal to
$p^2/2+p\log_2p+\BigO(p)$.
\end{proof}

\begin{remark}[Intrinsic endpoint cost]
The quadratic precision loss is not an artifact of the spline evaluator or of
bisection.  The two admissible inputs $0$ and $F(2^{-p})$ have outputs separated
by $2^{-p}$.  Any uniform oracle algorithm that promises $p$ output bits must,
in the worst endpoint case, obtain enough input information to distinguish
those values.  The exact modulus quantifies that information barrier.
\end{remark}

\begin{corollary}[Growth of exact and explicit reciprocal moduli]\label{co:cor:d-growth}
Let
\[
 d_{\mathrm{opt}}(n)=\Ceiling{1/F(1/n)},\qquad
 d_*(n)=\Ceiling{1/F(2^{-r(n)})},
\]
and let $d_\Delta(n)$ denote the explicit denominator $d(n)$ of
\cref{co:thm:reciprocal-modulus}.  Then each of
$d_{\mathrm{opt}},d_*$, and $d_\Delta$ has
\begin{equation}\label{co:eq:d-growth}
 \log_2 d_\bullet(n)
 =\frac12(\log_2 n)^2
  +(\log_2n)\log_2\log_2n
  +\BigO(\log n).
\end{equation}
Equivalently,
\begin{equation}\label{co:eq:d-quasipoly}
 d_\bullet(n)=n^{\frac12\log_2n+\log_2\log_2n+\BigO(1)}.
\end{equation}
Here $d_{\mathrm{opt}}$ is the mathematically least denominator for the
$1/n$ target, $d_*$ is recursive and exact-dyadic, and $d_\Delta$ is the
closed elementary witness; no computability claim is made here for
$d_{\mathrm{opt}}$.
\end{corollary}

\begin{proof}
Put $p=r(n)$.  Minimality of $p$ gives
\[
 2^{-p}<\frac1n\le2^{-(p-1)},
 \qquad p=\log_2n+\BigO(1).
\]
Since $F$ is increasing, with
$\Lambda(q)=\log_2(1/F(2^{-q}))$ this implies
\[
 \Lambda(p-1)\le \log_2\frac1{F(1/n)}<\Lambda(p).
\]
By \cref{co:thm:mu-asymptotic}, and because
$\mu(q)=\Ceiling{\Lambda(q)}$,
\[
 \Lambda(q)=\frac{q^2}{2}+q\log_2q+\BigO(q).
\]
The preceding sandwich therefore gives the same formula with $q=p$ for
$\log_2(1/F(1/n))$.  For every $x\ge1$,
$x\le\Ceiling{x}<2x$, so applying the outer ceiling changes its base-$2$
logarithm by at most one.  This proves \eqref{co:eq:d-growth} for
$d_{\mathrm{opt}}$.

The same ceiling observation applied directly to $F(2^{-p})$ proves it for
$d_*$.  Finally the closed formula for $d_\Delta$ gives
\[
 \log_2d_\Delta(n)
 =\frac{p(p+1)}2+(p+1)\log_2(p+1)
 =\frac{p^2}{2}+p\log_2p+\BigO(p).
\]
Substitute $p=\log_2n+O(1)$ in all three cases.  Exponentiating the resulting
identity proves \eqref{co:eq:d-quasipoly}.
\end{proof}

\section{Sequential computability by tolerant bisection}
\label{co:sec:sequential}

\subsection{Why ordinary bisection is insufficient}

Suppose $y$ is given only by a computable name.  At a dyadic midpoint $c$,
one can approximate both $F(c)$ and $y$, but in general no algorithm can
decide whether the two represented reals are exactly equal.  An implementation
that loops until it proves either $F(c)<y$ or $F(c)>y$ can therefore fail to
terminate when $y=F(c)$.

The inverse modulus provides the missing third branch.  If the approximations
are too close to certify a sign, then $F(c)$ is close to $y$; effective
injectivity says that $c$ is close to $G(y)$, so the algorithm may return it.

\subsection{The algorithm on a computable sequence}

Let $(y_i)$ be a computable sequence in $[0,1]$.  Fix a naming algorithm
$q(i,s)\in\RationalNumbers$ with
\begin{equation}\label{co:eq:input-name}
 \AbsoluteValue{q(i,s)-y_i}\le2^{-s}.
\end{equation}
We describe an algorithm which, from $(i,p)$, returns a dyadic $z(i,p)$ with
\begin{equation}\label{co:eq:output-name-goal}
 \AbsoluteValue{z(i,p)-G(y_i)}<2^{-p}.
\end{equation}

\begin{algorithm}[Tolerant certified bisection]\label{co:alg:tolerant-bisection}
Given an index $i$ and output precision $p$:
\begin{enumerate}
\item Set
\[
 r=p+2,
 \qquad h=2^{-r},
 \qquad \delta=\Delta_r.
\]
Choose integers $s,N$ effectively so that
\begin{equation}\label{co:eq:oracle-precisions}
 2^{-s}\le\delta/8,
 \qquad 2^{-N}\le\delta/8.
\end{equation}
Compute the rational $\widetilde y=q(i,s)$.
\item Initialize the certified bracket $[a,b]=[0,1]$.
\item Repeat at most $p+2$ times:
  \begin{enumerate}
  \item Put $c=(a+b)/2$ and evaluate the exact rational $S_N(c)$ by
        \eqref{co:eq:centered-spline-formula}.
  \item Form the rational approximate difference
        \begin{equation}\label{co:eq:D-tilde}
         \widetilde D=S_N(c)-\widetilde y.
        \end{equation}
  \item If $\widetilde D>\delta/2$, replace $b$ by $c$.
  \item If $\widetilde D<-\delta/2$, replace $a$ by $c$.
  \item Otherwise return $c$.
  \end{enumerate}
\item If no midpoint was returned, return $(a+b)/2$ after the final update.
\end{enumerate}
\end{algorithm}

Every comparison in \cref{co:alg:tolerant-bisection} is between exact rationals.
The order $N$ can be extremely large, but it is obtained by a finite integer
calculation and the spline sum is finite.  Hence the procedure is an algorithm
in the computability-theoretic sense.

\subsection{Correctness of every branch}

\begin{lemma}[Difference certificate]\label{co:lem:difference-certificate}
At every midpoint $c$ in \cref{co:alg:tolerant-bisection},
\begin{equation}\label{co:eq:difference-certificate}
 \AbsoluteValue{\widetilde D-(F(c)-y_i)}\le\delta/4.
\end{equation}
\end{lemma}

\begin{proof}
By \cref{co:thm:centered-error} and \eqref{co:eq:oracle-precisions},
$|S_N(c)-F(c)|\le\delta/8$.  By \eqref{co:eq:input-name} and the same choice of
$s$, $|\widetilde y-y_i|\le\delta/8$.  The triangle inequality gives
\eqref{co:eq:difference-certificate}.
\end{proof}

\begin{lemma}[Safe left/right updates]\label{co:lem:safe-updates}
If $\widetilde D>\delta/2$, then $G(y_i)<c$.  If
$\widetilde D<-\delta/2$, then $c<G(y_i)$.
\end{lemma}

\begin{proof}
In the first case, \cref{co:lem:difference-certificate} gives
\[
 F(c)-y_i\ge\widetilde D-\delta/4>\delta/4>0.
\]
Since $F$ is strictly increasing, $c>G(y_i)$.  The second case is symmetric.
Thus every update preserves the invariant
\begin{equation}\label{co:eq:bracket-invariant}
 G(y_i)\in[a,b].
\end{equation}
\end{proof}

\begin{lemma}[The inconclusive branch is a success branch]\label{co:lem:near-branch}
If $|\widetilde D|\le\delta/2$, then
\begin{equation}\label{co:eq:near-output}
 \AbsoluteValue{c-G(y_i)}<h=2^{-r}<2^{-p}.
\end{equation}
\end{lemma}

\begin{proof}
By \cref{co:lem:difference-certificate},
\[
 |F(c)-y_i|\le|\widetilde D|+\delta/4
 \le3\delta/4<\delta.
\]
The choice $\delta=\Delta_r$ and \cref{co:thm:closed-dyadic-modulus}, applied to
the input pair $F(c),y_i$, give
\[
 |G(F(c))-G(y_i)|<2^{-r}.
\]
Because $c\in[0,1]$, $G(F(c))=c$.  Finally, $r=p+2$.
\end{proof}

\begin{theorem}[Correctness and uniformity of tolerant bisection]
\label{co:thm:bisection-correct}
Algorithm \ref{co:alg:tolerant-bisection} terminates and returns a dyadic
$z(i,p)$ satisfying \eqref{co:eq:output-name-goal}.  The map $(i,p)\mapsto z(i,p)$
is computable whenever the name $(i,s)\mapsto q(i,s)$ is computable.
\end{theorem}

\begin{proof}
The loop has a fixed finite bound.  If it returns through the inconclusive
branch, \cref{co:lem:near-branch} proves the error estimate.  Otherwise
\cref{co:lem:safe-updates} preserves the bracket invariant through $p+2$
bisections.  Its final width is $2^{-(p+2)}$, so the final midpoint is at
distance at most $2^{-(p+3)}<2^{-p}$ from $G(y_i)$.

All quantities $r,h,\delta,s,N$ are computable from $p$.  Midpoints are
dyadic; $S_N(c)$ is an exact finite rational expression; and every branch is
a rational comparison.  The only access to the sequence is the single call to
its naming algorithm.  Therefore the construction is uniform in $i$ and $p$.
\end{proof}

\begin{corollary}[Sequential computability on the unit interval]
\label{co:cor:G-sequential}
The inverse $G:[0,1]\to[0,1]$ is sequentially computable.
\end{corollary}

\begin{proof}
Given any computable sequence $(y_i)$ in $[0,1]$, the output of
\cref{co:alg:tolerant-bisection} is a computable fast dyadic name for the sequence
$(G(y_i))$.
\end{proof}

\subsection{Clamping and the totalized inverse}

\begin{lemma}[Computability of clamping]\label{co:lem:clamp-computable}
If $(y_i)$ is a computable real sequence, then
$(\Clamp(y_i))$ is a computable sequence.
\end{lemma}

\begin{proof}
Clamping is $1$-Lipschitz.  If $q(i,p)$ approximates $y_i$ within $2^{-p}$,
then the exactly rational quantity $\Clamp(q(i,p))$ approximates
$\Clamp(y_i)$ within the same error.  The rational min/max operations are
decidable and computable.
\end{proof}

\begin{corollary}[Sequential computability of the totalized inverse]
\label{co:cor:GT-sequential}
The totalized inverse $\InvT:\RealNumbers\to\RealNumbers$ is sequentially computable.
\end{corollary}

\begin{proof}
Given a computable sequence $(y_i)$, first compute the clamped sequence by
\cref{co:lem:clamp-computable}, then apply \cref{co:cor:G-sequential}.  By
\eqref{co:eq:totalized-inverse}, the result is $(\InvT(y_i))$.
\end{proof}

\subsection{Main-theorem dependency checkpoint}

At this point every result cited in the proof of \cref{co:thm:main} has been
established: tolerant bisection supplies one uniform realizer transformer,
clamping totalizes it, and the closed dyadic threshold supplies the recursive
reciprocal modulus.  Thus the early proof is now entirely discharged rather
than relying on a compactness or choice argument.

\subsection{An abstract inverse-computability principle}

The Fabius proof illustrates a reusable theorem.  Stating it clarifies which
properties are genuinely needed and which belong only to this example.

\begin{theorem}[Effective inversion on a computable compact interval]
\label{co:thm:abstract-inversion}
Let $f:[0,1]\to[0,1]$ be a strictly increasing homeomorphism.  Suppose:
\begin{enumerate}
\item there is one algorithm which, given a dyadic rational $c\in[0,1]$ and
      $s\in\NaturalNumbers$, returns a rational within $2^{-s}$ of $f(c)$;
\item there is a computable positive rational sequence $(\alpha_p)$ such that
      for all $0\le x\le1-2^{-p}$,
      \begin{equation}\label{co:eq:abstract-gap}
       f(x+2^{-p})-f(x)\ge\alpha_p.
      \end{equation}
\end{enumerate}
Then $f^{-1}$ is sequentially computable and effectively uniformly continuous.
A tolerant-bisection algorithm is obtained by replacing $\Delta_r$ with
$\alpha_r$ in \cref{co:alg:tolerant-bisection}.
\end{theorem}

\begin{proof}
Condition \eqref{co:eq:abstract-gap} implies
$|u-v|<\alpha_p\Rightarrow|f^{-1}(u)-f^{-1}(v)|<2^{-p}$ by the same
contrapositive argument as \cref{co:cor:effective-injectivity}: if two inverse
values were at least $2^{-p}$ apart, monotonicity and
\eqref{co:eq:abstract-gap} would force their inputs to be at least
$\alpha_p$ apart.

This implication gives the uniform realizer.  On input precision $p$, run the
tolerant bisection of \cref{co:alg:tolerant-bisection} for $p+2$ steps, replacing
$\Delta_{p+2}$ by the positive rational $\alpha_{p+2}$.  Use hypothesis~1 and
the input name to approximate both sides within $\alpha_{p+2}/8$.  A certified
sign preserves the inverse bracket; in the inconclusive branch their true
difference is below $3\alpha_{p+2}/4$, so the preceding implication certifies
output error below $2^{-(p+2)}$.  If that branch never occurs, the final
bracket has width $2^{-(p+2)}$.  All loop bounds, precisions, and comparisons
are effective, so this is one uniform name transformer and proves sequential
computability.

For the reciprocal convention, set $d(0)=1$.  For $n\ge1$, compute any
$p(n)\ge1$ with $2^{-p(n)}<1/n$ and define
\[
 d(n):=\Floor{\frac1{\alpha_{p(n)}}}+1.
\]
Because $\alpha_{p(n)}$ is a positive computable rational, $d$ is a positive
recursive integer function and $d(n)>1/\alpha_{p(n)}$, hence
$1/d(n)<\alpha_{p(n)}$.  Therefore
\[
 |u-v|<\frac1{d(n)}
 \Longrightarrow |f^{-1}(u)-f^{-1}(v)|<2^{-p(n)}<\frac1n.
\]
This is effective uniform continuity in the subset-domain sense.
\end{proof}

For the Fabius function, \cref{co:thm:least-mass} identifies the best possible
choice $\alpha_p=F(2^{-p})$; its minimizing intervals are exactly the left and
right endpoint intervals.  The box event in \cref{co:lem:box-event} supplies a
closed lower substitute requiring no invocation of the exact dyadic evaluator.

\section{Analytic consequences of the exact modulus}

\subsection{Uniform continuity without any H\"older exponent}

Effective uniform continuity is sometimes mentally associated with a power-law
modulus.  The inverse Fabius function is a clean counterexample to that
intuition.

\begin{theorem}[No positive H\"older exponent at the endpoints]
\label{co:thm:no-holder}
For every $\alpha>0$ and every $C>0$, there are arbitrarily small $y>0$ such
that
\begin{equation}\label{co:eq:no-holder-zero}
 G(y)>Cy^\alpha.
\end{equation}
By reflection, there are arbitrarily small $1-y>0$ such that
\begin{equation}\label{co:eq:no-holder-one}
 1-G(y)>C(1-y)^\alpha.
\end{equation}
Thus neither $G$ nor $\InvT$ is H\"older continuous of any positive order at
$0$ or $1$, despite being effectively uniformly continuous.
\end{theorem}

\begin{proof}
Choose an integer $q$ with $q\alpha>1$.  By
\cref{co:lem:flatness-upper}, for all sufficiently small $x>0$,
$F(x)\le C_qx^q$ with a fixed $C_q$.  If $G(y)\le Cy^\alpha$ held near zero,
substitution $y=F(x)$ would give
\[
 x=G(F(x))\le C F(x)^\alpha
 \le C C_q^\alpha x^{q\alpha}.
\]
After division by $x$, the right side tends to zero, a contradiction.  The
right endpoint follows from $G(1-y)=1-G(y)$.
\end{proof}

\begin{corollary}[The exact modulus is sub-H\"older]\label{co:cor:sub-holder}
For every $\alpha>0$,
\begin{equation}\label{co:eq:sub-holder}
 \frac{\omega_G(\eta)}{\eta^\alpha}
 =\frac{G(\eta)}{\eta^\alpha}\longrightarrow+\infty
 \qquad(\eta\downarrow0).
\end{equation}
\end{corollary}

\begin{proof}
Choose an integer $q$ with $q\alpha>1$ and write $x=G(\eta)$, so that
$\eta=F(x)$ and $x\downarrow0$ with $\eta\downarrow0$.  By
\cref{co:lem:flatness-upper}, for all sufficiently small $x$,
$F(x)\le C_qx^q$.  Hence
\[
 \frac{G(\eta)}{\eta^\alpha}
 =\frac{x}{F(x)^\alpha}
 \ge C_q^{-\alpha}x^{1-q\alpha}\longrightarrow+\infty.
\]
The equality with $\omega_G$ is \cref{co:thm:exact-inverse-modulus}.
\end{proof}

The endpoint growth used here is also formalized, independently of the new
exact-modulus module, in \path{FabiusInverseAsymptotic.lean}.  The declarations
\path{Fabius.rpow_isLittleO_fabiusInv_at_zero_right} and
\path{Fabius.tendsto_fabiusInv_div_rpow_atTop_at_zero_right} give the
sub-H\"older divergence, while
\path{Fabius.fabiusInv_not_isBigO_rpow_at_zero_right} and
\path{Fabius.not_exists_fabiusInv_le_const_mul_rpow_near_zero} give the
corresponding impossibility of a positive power-law upper bound.  These are
analytic inputs already present in the repository, not consequences claimed
to have been newly proved by \path{InverseModulus.lean}.

\subsection{Endpoint asymptotics become modulus asymptotics}

The live formal corpus proves the leading inverse equivalent through
\lean{fabiusInv_isEquivalent_fabiusInverseAsymptoticMain} in
\path{FabiusInverseAsymptotic.lean}.  The companion frontier reports give
stronger paper-level inverse expansions\ \cite{is:p1:ProveItPrimary}.
Let
\begin{equation}\label{co:eq:T-L-rho}
 T=\log(1/\eta),
 \qquad L=\log2,
 \qquad \rho=\sqrt{\frac{2T}{L}}.
\end{equation}
Suppressing the paper-level first log-periodic correction, the formalized leading
equivalent is
\begin{equation}\label{co:eq:G-leading}
 G(\eta)\sim
 G_0(\eta):=
 \frac{\sqrt T}{\EulerE\sqrt L}
 \exp\!\bigl(-\sqrt{2LT}\bigr)
 \qquad(\eta\downarrow0).
\end{equation}
More precisely, the all-orders inversion theorem
\cref{ao:thm:all-orders} gives the relative expansion
\begin{equation}\label{co:eq:G-periodic-correction}
 \frac{G(\eta)}{G_0(\eta)}
 =1+
 \frac{\frac12\log\rho+\kappa_0-\Psi(\theta)}{\rho}
 +\BigO\!\left(\frac{(\log\rho)^2}{\rho^2}\right),
\end{equation}
where $\Psi$ is a nonconstant smooth one-periodic Gamma--zeta fluctuation and
$\theta:=\rho+\beta$ is the leading saddle phase.  It is not the exact
lower-Lambert phase: the latter is
\[
 \lambda(G(\eta))=-\frac1L\LambertW_{-1}(-LG(\eta))
 =\theta+\BigO\!\left(\frac{\log\rho}{\rho}\right).
\]
Thus replacing the exact phase by $\theta$ is justified at the displayed
algebraic order, but the distinction matters for exponentially small terms.

\begin{corollary}[Paper-level transfer of the frontier inverse expansion]
\label{co:cor:sharp-asymptotic-modulus}
At the paper level, the expressions \eqref{co:eq:G-leading} and
\eqref{co:eq:G-periodic-correction} transfer, without any further inversion, to
the leading and first-correction asymptotics of the exact continuity modulus
of $G$ and $\InvT$:
\begin{equation}\label{co:eq:asymptotic-modulus}
 \omega_G(\eta)=\omega_{\InvT}(\eta)=G(\eta)
 \qquad(0\le\eta\le1).
\end{equation}
Likewise, every valid all-orders endpoint expansion for the inverse transfers
verbatim to an all-orders modulus-of-continuity expansion.  Lean currently
proves the leading equivalent and an implicit expansion pulled back to exact
inverse coordinates, but not the displayed explicit periodic correction or
an explicit all-orders reversion.
\end{corollary}

\begin{proof}
This is exactly \cref{co:thm:exact-inverse-modulus}.
\end{proof}

This observation closes a conceptual loop in the corpus.  The lower-Lambert
phase and Bell-polynomial inverse transfer were developed to approximate
quantiles.  The geometric theorem shows that the same quantiles are the
optimal global sensitivity profile of the inverse map.

\subsection{Derivative blow-up versus computability}

On $0<y<1$, the repository proves
\begin{equation}\label{co:eq:inverse-derivative}
 G'(y)=\frac{1}{F'(G(y))}
 =\frac{1}{2\RvachevUp(2G(y)-1)}>0.
\end{equation}
The derivative tends to $+\infty$ at both endpoints.  There is therefore no
global Lipschitz constant and no endpoint condition number bounded in the
ordinary differential sense.  Nonetheless the exact nonlinear modulus
$G(\eta)$ is computable, and tolerant bisection uses precisely that nonlinear
conditioning.  This is why derivative-based numerical intuition alone is not
a substitute for a computable-analysis proof.

% Source fragment: Inverse computability; prefix co.
\section{Further deductions, conjectures, and research directions}
\label{co:sec:frontiers}

\subsection{A certified practical inverse evaluator}

The proof algorithm is intentionally elementary rather than efficient.  A
practical certified implementation should retain its logical shell while
replacing the forward oracle by a portfolio of enclosures:

\begin{itemize}
\item exact dyadic evaluation and terminating local Taylor formulae near
      dyadic anchors;
\item centered Thue--Morse splines at moderate depth;
\item quarter-base $q$-Richardson acceleration of finite-prefix errors;
\item Fourier/sinc evaluation in regions where frequency-space decay is
      favorable;
\item lower-Lambert endpoint enclosures when the requested target is extremely
      close to $0$ or $1$.
\end{itemize}

Each backend must return a rational interval containing $F(c)$, not merely a
floating-point approximation.  The bisection controller only asks whether
that interval lies wholly above or below the target interval; overlapping
intervals trigger refinement or the modulus-based acceptance branch.  Thus
backend selection is a performance concern, not part of the correctness
argument.

\begin{problem}[Near-optimal certified evaluation]
\label{comp:prob:near-optimal-evaluation}
Construct and formally verify an inverse evaluator whose oracle-bit
consumption is $\mu(p)+\BigO(\log p)$ for $p$ output bits at the worst endpoint,
and whose arithmetic running time is quasi-linear or softly quadratic in the
size of the exact integers produced.  Separate information-theoretic oracle
precision from the cost of evaluating $F$ to that precision.
\end{problem}

\subsection{Sharp integer moduli and the log-periodic phase}

The real-valued optimal input-bit threshold is
\[
 \Lambda(p):=\log_2\frac1{F(2^{-p})},
 \qquad \mu(p)=\Ceiling{\Lambda(p)}.
\]
The all-orders endpoint theory expresses $\Lambda(p)$ in a lower-Lambert
phase with a nonconstant one-periodic Gamma--zeta contribution and a hierarchy
of inverse powers.  Passing to the integer ceiling introduces a second,
sawtooth nonlinearity.  It is therefore natural to study the arithmetic
interaction between the smooth periodic saddle phase and fractional parts of
its growing nonperiodic carrier.

For $p\ge1$ define the exact lower-Lambert phase
\[
 \lambda_p:=-\frac1L\LambertW_{-1}(-L2^{-p}),
 \qquad 2^{-p}=\lambda_p e^{-L\lambda_p}.
\]
The first valid forward remainder gives the precise reduction
\begin{equation}\label{co:eq:Lambda-lower-Lambert}
 \Lambda(p)=\frac{\lambda_p^2}{2}-\beta\lambda_p
 +\frac{\log\lambda_p}{2L}-\frac{\Csharp}{L}
 -\frac{\Psi(\lambda_p)}{L}+\BigO(\lambda_p^{-1}).
\end{equation}
Indeed, substitute $x=2^{-p}$ and its exact phase $\lambda_p$ into
\eqref{ao:eq:forward-all-orders} with $N=1$, negate, and divide by
$L=\log2$; this is exactly the definition of $\Lambda(p)$.
Put
\[
 A(p):=\frac{\lambda_p^2}{2}-\beta\lambda_p
 +\frac{\log\lambda_p}{2L}-\frac{\Csharp}{L},
 \qquad R(p):=\mu(p)-A(p).
\]
Equation \eqref{co:eq:Lambda-lower-Lambert} and
$0\le\Ceiling{\Lambda(p)}-\Lambda(p)<1$ show that $R(p)$ is bounded.

\begin{conjecture}[Nontrivial fluctuation of the optimal integer modulus]
\label{co:conj:integer-modulus}
The bounded sequence $(R(p))_{p\ge1}$ has at least two accumulation points and
is not eventually periodic.
\end{conjecture}

The conjecture isolates a definite sequence and makes no unsupported claim
that its accumulation set is a simple sum of two independent sets.  A stronger
description or equidistribution theorem would require joint information about
$\lambda_p\bmod1$ and the fractional parts of $\Lambda(p)$, which the present
corpus does not supply.  Exact values $F(2^{-p})\in\RationalNumbers$ make the
sequence amenable to arbitrary-precision experimentation without numerical
ambiguity.

\subsection{Base-\texorpdfstring{$b$}{b} atomic distributions}

Replace the dyadic random series by
\[
 X_b=\sum_{k\ge1}b^{-k}V_k,
 \qquad b\ge2,
\]
with independent uniform coordinates.  The finite convolution is still
computable, but binary Thue--Morse signs are replaced by base-$b$ subset or
digit combinatorics, and the support is $[0,1/(b-1)]$ before normalization.
The least-interval-mass theorem continues to follow whenever the normalized
density is symmetric and unimodal.  The same proof would then identify the
optimal inverse modulus with the quantile itself.

\begin{problem}[Uniform computability across the base]
\label{comp:prob:uniform-base}
Develop a theorem uniform in the integer parameter $b$: construct a single
algorithm that, given $b$, a name of $y$, and a precision $p$, computes the
normalized quantile of $X_b$.  Determine the leading dependence of the
optimal input-bit modulus on both $b$ and $p$, and identify the corresponding
base-$b$ log-periodic correction.
\end{problem}

\subsection{Certified spectral and polynomial representations}

The Legendre, Lagrange, and shifted-$\RvachevUp$ representations in the corpus are
mathematically exact, but a representation becomes an algorithm only after a
recursive uniform tail bound is attached.  This suggests three concrete
projects:

\begin{enumerate}
\item derive explicit Legendre coefficient envelopes strong enough to certify
      pointwise and uniform evaluation of $\RvachevUp$;
\item combine the finite shifted-$\RvachevUp$ representation of polynomials with
      interval arithmetic to obtain certified local surrogates for $F$;
\item compare physical-space spline, Legendre, and sinc-product evaluators by
      bit complexity, not merely by observed decimal convergence.
\end{enumerate}

A cross-certified implementation could evaluate the same point by two
independent representations and intersect the resulting enclosures.  This is
particularly attractive near transition regions where one representation's
natural error bound is pessimistic.

\subsection{Higher-dimensional monotone transport}

The one-dimensional theorem
\[
 \min_x \Probability\{X\in(x,x+h]\}=\Probability\{X\in(0,h]\}
\]
is a sharp anti-concentration lower bound tailored to a symmetric unimodal
law.  Because the law is atomless, the same identity holds with any choice of
open or closed endpoints.  In higher dimensions, inverse CDFs are replaced by
triangular Rosenblatt maps or optimal-transport maps.  Effective invertibility
would require computable lower bounds on masses of suitable boxes or sections.

\begin{problem}[Atomic transport computability]
\label{comp:prob:atomic-transport}
For products or coupled variants of Fabius-type laws, determine conditions
under which the triangular transport from Lebesgue measure has a computable
inverse with an explicit multivariate modulus.  Investigate whether total
positivity or variation-diminishing properties supply the needed lower
section-mass estimates.
\end{problem}

\subsection{Branchwise inversion of the signed global extension}

The signed global Fabius extension is not globally injective, so it has no
single inverse.  On every maximal monotonicity interval, however, it has a
branch inverse.  Local finiteness of its Thue--Morse translate expansion makes
forward evaluation computable.  The remaining issue is effective branch
separation near critical values.

\begin{problem}[Computable branch atlas]
\label{comp:prob:branch-atlas}
Construct an effective enumeration of monotonicity intervals of the signed
global extension, together with certified branch ranges and recursive inverse
moduli.  Determine whether the self-similar translation law permits one finite
branch template plus a Thue--Morse address.
\end{problem}
